\section{Channels outside the framework}\label{sec:22.3}

A dark-matter fermion of \SI{1}{MeV} absorbed by a $^{125}$Te nucleus in a TeO$_2$ crystal can turn a neutron into a proton and leave as an electron with \SI{626}{keV} of kinetic energy. The rate of that process carries a factor of 6.3 that no response of the crystal contains: the attraction that the daughter nucleus, of charge 53, exerts on the outgoing electron. The master formula \cref{eq:8.2-rate} holds for every process in which the dark-matter side can be traced out before the target, so that a tensor $\mathfrak{D}_{ab}(\vect{q}, \vect{v})$ that depends only on the dark-matter side multiplies a response $S_{ab}(\vect{q}, \omega)$ that depends only on the target. Where that separation fails, the process lies outside the framework of this book.

\subsection{The scope of the master formula}\label{sec:22.3-scope}

The derivation of \cref{eq:8.2-rate} in \cref{app:D.25} uses three properties of the amplitude besides the assumptions of \cref{sec:8.2-conditions}. The dark field acts through the rows of \cref{eq:8.1-rows}, one-body operators that conserve the number of particles of every species, once per amplitude (twice for the processes of \cref{sec:10.8}). Whatever leaves on the dark-matter side is a free plane wave. The matrix element is then a dark-matter factor that depends on $\vect{q}$ and $\vect{v}$ times $\langle f|O_a(\vect{q})|i\rangle$, and the sum over final states of the target closes into a two-point function at one momentum transfer. The target starts in a Gibbs state, and the observable is the excitation it is left in, read through a final-state class of \cref{sec:8.7}. A process that violates the first property needs a row that \cref{eq:8.1-rows} lacks. One that violates the second has no $\mathfrak{D}_{ab}$ at all, because the outgoing particle acts on the target. One that violates the third can have a well-defined rate that no detector reads as an excitation spectrum.

Two of the assumptions of \cref{sec:8.2-conditions} also bound the framework. With couplings near $4\pi$ and a mediator lighter than about $20A\mu_{\chi T}$, with $\mu_{\chi T}$ the reduced mass of dark matter and nucleus, an attractive potential binds and the first Born approximation fails at order one (\cref{sec:8.2-conditions}). A target whose density operator has coherences between energy eigenstates, or one driven during the exposure, has no response of the form \cref{eq:4.2-Sab}. A last class lies inside the framework: the rate is defined, but no calculation of $S_{ab}$ for that kind of target yet obeys the exact relations of \cref{sec:8.3}. \Cref{tab:22-scope} lists the processes that break one of the three properties, the two that the assumptions exclude, and the one inside the framework that has no conserving response yet.

{\scriptsize\setlength{\tabcolsep}{3pt}\setlength{\LTcapwidth}{\textwidth}
\begin{longtable}{@{}>{\raggedright\arraybackslash}p{0.17\linewidth}>{\raggedright\arraybackslash}p{0.27\linewidth}>{\raggedright\arraybackslash}p{0.38\linewidth}>{\raggedright\arraybackslash}p{0.12\linewidth}@{}}
  \caption{Processes outside the framework of the master formula, and one inside it without a computed response. Second column: why the rate is not $\sum_{ab}\mathfrak{D}_{ab}S_{ab}$ of \cref{eq:8.2-rate}, or, for the last row, why no computed response is available. Third column: the calculation or measurement that a prediction needs. Last column: the section that states the process.}
  \label{tab:22-scope}\\
  \toprule
  Process & Reason & What a prediction needs & Section \\
  \midrule
  \endfirsthead
  \toprule
  Process & Reason & What a prediction needs & Section \\
  \midrule
  \endhead
  \bottomrule
  \endfoot
  Charged-current absorption, $\chi + (A, Z) \to (A, Z{+}1) + e^-$ & Species-changing nuclear row; the outgoing electron moves in the field of the final state of the target, so the dark-matter side is an operator on the target & Shell-model Fermi and Gamow--Teller strengths per isotope, with quenched axial coupling; Dirac--Coulomb electron waves for the daughter charge & \cref{sec:22.3-cc} \\
  \midrule
  Radiative decay of an excited dark state & No target operator in the amplitude; the medium enters through the photon propagator only & Abundance of the excited state at the detector; local field at the decaying particle; absorption of the photon & \cref{sec:10.9-decay,sec:22.3-tracks} \\
  \midrule
  Damage tracks of recoiling heavy ions & Born parameter far above one; the lattice is damaged by the earlier collisions; the readout is a defect configuration & Transport of the ion with the stopping as input, and the map from the track to the defects that are read & \cref{sec:19.3-bethe,sec:22.3-tracks} \\
  \midrule
  Strong dark-matter--target coupling & First Born approximation fails; bound states and resonances & Partial-wave solution of the scattering & \cref{sec:8.2-conditions} \\
  \midrule
  Target with coherences, or driven during the exposure & No Gibbs or energy-diagonal initial state & Rate for a density operator with coherences between energy eigenstates & \cref{sec:8.2-conditions} \\
  \midrule
  Charge-density-wave phasons and amplitudons & Inside the framework; the response needs a conserving kernel on the mean field of the distorted crystal & Conserving time-dependent mean-field response; Goldstone and $f$-sum checks & \cref{sec:22.3-cdw} \\
\end{longtable}}

\subsection{Charged-current absorption}\label{sec:22.3-cc}

The fermionic absorption of \cref{sec:8.1-absorption} has a charged partner. An operator that couples the dark-matter fermion to an electron and a neutron to a proton,
\begin{equation}
  \frac{1}{\Lambda_{\mathrm{F}}^2}\,\bigl[\bar p\,\gamma_\mu(1 + \lambda_A\gamma^5)\,n\bigr]\bigl[\bar e\,\gamma^\mu P_R\,\chi\bigr] + \text{h.c.} ,
  \label{eq:22.3-cc}
\end{equation}
with $p$, $n$, $e$ and $\chi$ the proton, neutron, electron and dark-matter fields, $\Lambda_{\mathrm{F}}$ the scale of the operator as for the neutral absorption operators of \cref{sec:8.1-absorption}, $P_R = (1 + \gamma^5)/2$ and the nucleon axial coupling $\lambda_A = 1.27$ in the sign convention of the reference, absorbs $\chi$ on a neutron of a nucleus and emits an electron, $\chi + (A, Z) \to (A, Z+1)^{(*)} + e^-$~\cite[Eq.~(4.1)]{dror2020absorption}; its pure vector form is Eq.~(6) of Ref.~\cite{dror2020direct}. For a parent that is stable against $\beta$ decay the process has a threshold, $\mchi > m_{\mathrm{th}}^{\mathrm{CC}} = M_{A,Z+1}^{(*)} + m_e - M_{A,Z}$ with nuclear masses, the daughter possibly excited~\cite[Eq.~(4.2)]{dror2020absorption}. Above it the electron takes the kinetic energy $E_{\mathrm{ej}} = \mchi - m_{\mathrm{th}}^{\mathrm{CC}}$, up to a nuclear recoil of order $\mchi^2/2M$. The ground state of $^{125}$Te has spin $\frac12^+$, and the ground state of $^{125}$I, $\frac52^+$, cannot be reached from it by an allowed transition; the lowest allowed threshold, to an excited state of $^{125}$I reached by a Gamow--Teller transition, is \SI{374}{keV}~\cite[Table~1]{dror2020absorption}.

The nuclear side is a row that \cref{eq:8.1-rows} does not have. At momentum transfers small against the inverse nuclear radius, the time component of the vector current reduces to the isospin raising operator $\tau^+_j$, which turns nucleon $j$ from a neutron into a proton. The spatial axial current reduces to $\sigma^\beta_j\tau^+_j$, as the density and the spin density of \cref{app:D.24} do for the neutral currents. Summed over the nuclei $I$ at positions $\vect{X}_I$,
\begin{equation}
  O^{\mathrm{CC}}_{\mathrm{F}}(\vect{q}) = \sum_I\ee^{\ii\vect{q}\cdot\vect{X}_I}\sum_{j\in I}\tau^+_j , \qquad
  O^{\mathrm{CC}}_{\mathrm{GT},\beta}(\vect{q}) = \sum_I\ee^{\ii\vect{q}\cdot\vect{X}_I}\sum_{j\in I}\sigma^\beta_j\tau^+_j ,
  \label{eq:22.3-rows}
\end{equation}
the Fermi and Gamow--Teller rows. Every row of \cref{eq:8.1-rows} conserves the numbers of protons and neutrons; these change the element. Adding these rows, and nothing else, would keep the framework, with a response $S^{\mathrm{CC}}_{ab}(\vect{q}, \omega)$ of the rows in \cref{eq:22.3-rows}.

The electron prevents that. It is charged, and it leaves through the Coulomb field of the daughter nucleus and of the atomic electrons around it. Its final state is a scattering state $\psi^{(-)}_{\vect{p}_{\mathrm{ej}}}(\vect{r})$ of that field, not a plane wave, and for a dark-matter particle of momentum $\vect{p}$ the amplitude is an overlap $\int\dd^3r\,\psi^{(-)*}_{\vect{p}_{\mathrm{ej}}}(\vect{r})\,\ee^{\ii\vect{p}\cdot\vect{r}}\langle f|\rho^{\mathrm{CC}}(\vect{r})|i\rangle$ with the density $\rho^{\mathrm{CC}}(\vect{r})$ whose Fourier components are the rows in \cref{eq:22.3-rows}. Its square involves the transition density at two points, $\langle i|\rho^{\mathrm{CC}\dagger}(\vect{r}')|f\rangle\langle f|\rho^{\mathrm{CC}}(\vect{r})|i\rangle$, weighted by a product of electron waves, so it is not a response at one momentum transfer. The field that distorts the wave, moreover, belongs to the final state $|f\rangle$: the charge $Z + 1$ of the daughter and the atomic configuration left behind when the nuclear charge changes. The dark-matter side therefore acts on the target, and no tensor $\mathfrak{D}_{ab}$ exists. The exact rate is the golden rule with the electron in the final state of the combined system, $\sum_f|\langle f; e^{(-)}_{\vect{p}_{\mathrm{ej}}}|H_{\mathrm{int}}|i; \chi\rangle|^2$. It is the spectral function of an operator that adds an electron to the target and turns a neutron into a proton: a response of the target with one more electron. A condensed-matter reader knows the structure from the x-ray edge, where the core hole left in the final state distorts the outgoing electron.

One limit restores a product. Let the electron wavelength be long against the nuclear radius, $p_{\mathrm{ej}}r_A \ll 1$ with $r_A = 1.2A^{1/3}\,\mathrm{fm}$. The wave can then be evaluated at the nucleus. Its squared amplitude relative to a plane wave is the relativistic Fermi function $F_{\mathrm{Fermi}}(Z_{\mathrm{d}}, E_{\mathrm{ej}})$ of \cref{eq:21.7-fermi}, the point-charge function of $\beta$ decay with its traditional prefactor 4~\cite[Eq.~(6)]{hayen2018high}, evaluated at the electron's kinetic energy $E_{\mathrm{ej}}$ with $Z_{\mathrm{d}} = Z + 1$ the charge of the daughter and with the Sommerfeld parameter $y_{\mathrm{S}} = \alpha Z_{\mathrm{d}}E_e/p_{\mathrm{ej}}$ of \cref{eq:21.7-beta}, where $E_e = m_e + E_{\mathrm{ej}}$ is the total energy of the electron. The finite size of the nucleus enters through a further factor $L_0$ that differs from one by a quantity of order $\alpha Z_{\mathrm{d}}E_er_A$. In this allowed approximation the rate per isotope is a sum over final nuclear states of $F_{\mathrm{Fermi}}(Z_{\mathrm{d}}, E_{\mathrm{ej}})$ times the electron phase space, proportional to $p_{\mathrm{ej}}E_e$, times the standard allowed combination of $\beta$ decay~\cite[Sect.~I]{hayen2018high}, the Fermi strength $B_{\mathrm{F}} = |\langle f|\sum_j\tau^+_j|i\rangle|^2$ plus $\lambda_A^2$ times the Gamow--Teller strength $B_{\mathrm{GT}} = \sum_\beta|\langle f|\sum_j\sigma^\beta_j\tau^+_j|i\rangle|^2$. The published projections multiply the nucleon-level amplitude by the same kind of function~\cite[Eq.~(4.9)]{dror2020absorption}. Their Eq.~(4.10) prints the prefactor $2(1 + \gamma_{\mathrm{d}})$, a convention that moves part of $L_0$ into the function, and $\Gamma(1 + \gamma_{\mathrm{d}})^2$ in the denominator where \cref{eq:21.7-fermi} has $\Gamma(1 + 2\gamma_{\mathrm{d}})^2$. The gamma-function factor alone is therefore larger by $\Gamma(1 + 2\gamma_{\mathrm{d}})^2/\Gamma(1 + \gamma_{\mathrm{d}})^2$, 3.22 at $Z_{\mathrm{d}} = 53$ and 4.0 as $Z_{\mathrm{d}} \to 0$, and the printed function does not reach the nonrelativistic limit.

\begin{workedexample}{The Fermi function of $^{125}$Te absorption}\label{we:22-cc}
Take $^{125}$Te in TeO$_2$, the Gamow--Teller threshold of \SI{374}{keV} and $\mchi = \SI{1}{MeV}$. The electron leaves with $E_{\mathrm{ej}} = \SI{626}{keV}$ and $p_{\mathrm{ej}} = \SI{1.016}{MeV}$, so $y_{\mathrm{S}} = 0.433$ for the daughter iodine, $Z_{\mathrm{d}} = 53$, and $p_{\mathrm{ej}}r_A = 0.031$ with $r_A = \SI{6.0}{fm}$. The allowed approximation holds. Its corrections in the nuclear operators are of order $(p_{\mathrm{ej}}r_A)^2 \approx 10^{-3}$, and the finite-size Coulomb correction $L_0$ is of order $\alpha Z_{\mathrm{d}}E_er_A = 0.013$, about \SI{1}{\percent}.

If the master formula applied, the rate would be the contraction of a dark-matter tensor, built from the operator \cref{eq:22.3-cc} and a plane-wave electron, with the nuclear response of \cref{eq:22.3-rows} at $q = p_{\mathrm{ej}}$. At these momenta the nuclear strength does not depend on $q$, so at fixed matrix element the rate would grow with the electron phase space, as $p_{\mathrm{ej}}E_e$. With the distorted wave it is $p_{\mathrm{ej}}E_eF_{\mathrm{Fermi}}$, and \cref{eq:21.7-fermi} gives $F_{\mathrm{Fermi}}(53, \SI{626}{keV}) = 6.34$.

The factor belongs to the final state. Evaluated with the charge of the parent, $Z_{\mathrm{d}} \to 52$, it is 6.06, \SI{4.3}{\percent} lower. A response sums over final states independently of the outgoing particle, and a dark-matter tensor depends on $\vect{q}$ and $\vect{v}$ alone, so neither can carry this dependence.

It also changes the threshold behavior. At $E_{\mathrm{ej}} = \SI{10}{keV}$, $F_{\mathrm{Fermi}} = 32.5$, and at \SI{1}{keV}, 102. The product $p_{\mathrm{ej}}E_eF_{\mathrm{Fermi}}$ tends to a constant, 0.23 of its value at \SI{626}{keV} for both energies, while $p_{\mathrm{ej}}E_e$ alone falls to 0.046 and 0.014 of it. The rate stays finite at the threshold, as the $\beta$ spectra of \cref{sec:21.7} do at their low-energy end, where a plane-wave electron would make it vanish as $p_{\mathrm{ej}}$.

At low energies the atomic electrons enter as well. The radius of the iodine $K$ shell, $a_0/Z_{\mathrm{d}}$, is \SI{998}{fm}, against $1/p_{\mathrm{ej}} = \SI{194}{fm}$ at \SI{626}{keV} and \SI{1942}{fm} at \SI{10}{keV}. Below about \SI{10}{keV} the wave is screened by electrons whose configuration, after the sudden change of the nuclear charge, is not stationary, and the distortion then depends on the atomic final state too. The factorized form is an approximation for one final state of the nucleus with the atomic shells as spectators, and it holds because $p_{\mathrm{ej}}r_A$ is small and $E_{\mathrm{ej}}$ is large; the general rate is a property of the target together with the electron.
\end{workedexample}

The nuclear strengths obey two sum rules. Write $T_\pm = \sum_j\tau^\pm_j$ for the summed raising and lowering operators of isospin, which obey the angular-momentum algebra with $T_z = (Z - N)/2$ in the convention in which a proton has $\tau_z = +1$, so that $T_+$ turns a neutron into a proton. For the Fermi strength we need the isospin limit, in which the ground state of the parent has total isospin $T_{\mathrm{iso}} = (A - 2Z)/2$, the smallest value its projection $T_z = -T_{\mathrm{iso}}$ allows. It is then the bottom state of its multiplet, and $T_-$, which would turn one more proton into a neutron, annihilates it. The summed strength is then a ladder norm,
\begin{equation*}
  \sum_fB_{\mathrm{F}} = \langle i|T_-T_+|i\rangle = T_{\mathrm{iso}}(T_{\mathrm{iso}} + 1) - T_{\mathrm{iso}}(T_{\mathrm{iso}} - 1) = 2T_{\mathrm{iso}} = A - 2Z ,
\end{equation*}
and $T_+$ takes the ground state into a single state, the isobaric analog of the parent. With $\sigma^\beta_j$ inserted no such state exists, but the commutator of the summed operators with their conjugates, summed over $\beta$, picks up $\sum_\beta\sigma^\beta\sigma^\beta = 3$, and
\begin{equation}
  \sum_fB_{\mathrm{F}} = A - 2Z , \qquad
  \sum_fB_{\mathrm{GT}}(n \to p) - \sum_fB_{\mathrm{GT}}(p \to n) = 3(A - 2Z) ,
  \label{eq:22.3-sumrules}
\end{equation}
the Fermi and Ikeda sum rules, where $A - 2Z$ is the neutron excess. The Ikeda rule is an operator identity and holds without isospin symmetry; the Fermi rule in this form, and the concentration of its strength in one state, need the isospin limit. The analog state lies above the parent by the Coulomb energy of the added proton less the neutron--proton mass difference. For $^{125}$Te that puts the threshold to reach it at about 13 to \SI{14}{MeV}: \SI{14.2}{MeV} for a uniformly charged sphere of radius $r_A$, and \SI{13.6}{MeV} once the exchange term of the Coulomb energy of that sphere is subtracted. Below that mass a vector charged current acts only through isospin-breaking admixtures, and the axial current carries the rate through Gamow--Teller transitions. The published projections count the nucleon-level rate once per neutron, $n_j \simeq A - Z = 73$ for $^{125}$Te~\cite[Eq.~(4.6)]{dror2020absorption}. The Fermi sum rule allows 21 in total, all of it in the analog state. On the Gamow--Teller side the same count gives $3(A - Z) = 219$ against the Ikeda difference of 63, before quenching. Shell-model Gamow--Teller strengths also exceed measured ones unless the axial coupling in the nucleus is reduced by a factor of about 0.75, or two-body currents are added~\cite{gysbers2019discrepancy}, which changes the weight $\lambda_A^2$ by a factor of about two.

A prediction for a crystal therefore needs three things for each isotope. The first is the nuclear strength: shell-model Fermi and Gamow--Teller strengths for every final state below the dark-matter mass, with the quenched axial coupling or the two-body currents, and their dependence on $q$ through the transition densities, with first-forbidden operators added where $p_{\mathrm{ej}}r_A$ reaches about 0.1 or the allowed matrix element vanishes. The computed strengths must exhaust \cref{eq:22.3-sumrules} in a space closed under $\sigma\tau$, which contains every spin--orbit partner of its orbitals. The second is the electron: Dirac--Coulomb waves of the daughter charge, with finite nuclear size and atomic screening, evaluated with the same nuclear radius as the strengths; at \SI{626}{keV} the parent charge would give a Fermi function \SI{4.3}{\percent} too low. The third is the final state that a detector reads: the electron deposits its energy as a fast primary of the kind that \cref{chap:18} follows, together with the de-excitation of the daughter. Capture of the electron into a bound orbital of the daughter leaves no free particle to carry the energy away, so it occurs only where $\mchi$ matches a transition, within the width of the final atomic state.

\subsection{Radiative decay and ion tracks}\label{sec:22.3-tracks}

An excited dark state $\chi_2$ that decays to $\chi_1$ and a photon inside the detector has an amplitude with no target operator in it (\cref{sec:10.9-decay}). The event rate is then the density of excited states at the detector times their decay rate times the volume, with no flux factor $\rho_\chi v/\mchi$ and no exposure in target mass, and the medium enters only through the photon density of states at $|\Delta_\chi|$. A prediction needs the abundance of $\chi_2$ at the detector, set by the early universe or by up-scattering in the Earth; the local field at the decaying particle, which the macroscopic $\varepsilon^{\mathrm{T}}$ of \cref{eq:10.9-decay} omits; and the absorption length of the photon in the detector, which the optical response gives once the photon exists. The part of the decay that hands energy directly to the electrons through a virtual photon lies inside the framework (\cref{sec:10.9-decay}).

A heavy ion leaves a damage track (\cref{sec:19.3-bethe}). Its rate of energy loss is the Bethe or friction stopping of \cref{eq:19.3-bethe,eq:19.3-Se}, a response quantity, but the track is not. A recoiling silicon ion at the speeds that dark matter gives it has a Born parameter $Z_1\alpha/v$, with $v$ the speed of the ion, of 117 to over a thousand, so no single scattering is perturbative. This is the parameter of the moving ion; the Born parameter $Z_\kappa\alpha/v_{\mathrm{F}}$ of \cref{sec:22.1}, set by the electrons it meets, measures the same failure at low speed. A track of $10^2$ to $10^4$ collisions runs through a lattice that the earlier collisions have already displaced, and the target is not in its Gibbs state for any collision after the first. What a track detector or a paleo-detector reads is the defect configuration that remains, not an excitation spectrum. A prediction needs a transport calculation of the ion through the damaged lattice, with the stopping as input and molecular dynamics for the collisions, and the map from the final configuration to the defects that the readout sees.

\subsection{Collective modes of a charge-density wave}\label{sec:22.3-cdw}

Below about \SI{260}{K} the quasi-one-dimensional crystal (TaSe$_4$)$_2$I forms a charge-density wave with a single-particle gap $2\Delta_{\mathrm{CDW}} \approx 250$ to \SI{300}{meV}. It carries a collective mode at \SI{0.23}{THz}, or \SI{0.951}{meV}, which Kim and collaborators identify as the massive phason of that wave~\cite{kim2023observation}. The density row of \cref{eq:8.1-rows} can excite modes that lie 263 to 315 times below the charge gap: with $v_{\mathrm{max}}$ the largest dark-matter speed in the laboratory (\cref{eq:1.2-vmax}), a dark-matter particle heavier than $2\omega/v_{\mathrm{max}}^2 = \SI{271}{eV}$ can deposit \SI{0.951}{meV}. Nothing in the master formula changes for such a target. What is missing is a response that obeys the exact relations of \cref{sec:8.3}.

In a charge-density wave the electron density is modulated, $\nd_e(\vect{r}) = \bar\nd_e + \nd_1\cos(\vect{Q}_{\mathrm{CDW}}\cdot\vect{r} + \varphi)$, with a wave vector $\vect{Q}_{\mathrm{CDW}}$ that is incommensurate with the lattice in (TaSe$_4$)$_2$I, and the lattice distorts with it. The order parameter is $\Delta_{\mathrm{CDW}}\ee^{\ii\varphi}$. The phason is an oscillation of $\varphi$, the amplitudon one of $|\Delta_{\mathrm{CDW}}|$. If the wave is incommensurate, a uniform shift of $\varphi$ costs no energy, and the phason is a Goldstone mode whose frequency vanishes linearly in its wave vector; impurities and commensurability pin it and open a small gap. A density probe couples to the phason in two places. Near $\vect{q} = \pm\vect{Q}_{\mathrm{CDW}} + \vect{G}$ it meets the modulation, $\delta\nd_e = -\nd_1\sin(\vect{Q}_{\mathrm{CDW}}\cdot\vect{r} + \varphi)\,\delta\varphi$, and reaching it there takes $\mchi \ge |\vect{Q}_{\mathrm{CDW}}|/2v_{\mathrm{max}}$, \SI{189}{keV} per keV of $|\vect{Q}_{\mathrm{CDW}}|$. At small $\vect{q}$ a gradient of $\varphi$ along the chains moves the condensed electrons with the wave and leaves a charge density proportional to that gradient. That charge feels the long-range Coulomb interaction, which lifts the longitudinal phason, whose wave vector has a component along $\vect{Q}_{\mathrm{CDW}}$, from zero to the plasma frequency of the condensed electrons. A phason whose wave vector is perpendicular to the chains carries no charge and stays acoustic, or pinned, so the response depends on the direction of $\vect{q}$. Lee, Rice and Anderson, and Lee and Fukuyama, showed this, as quoted by Kim and collaborators~\cite{kim2023observation}. The condensate plasma frequency is as low as a meV because a moving charge-density wave drags its lattice distortion with it, which gives the condensed electrons a large effective mass~\cite{lee1974conductivity}. The superconductor of \cref{sec:15.1} shows the same mechanism, with its phase mode pushed to the plasmon.

The response that captures these modes has the structure of \cref{sec:15.1}. On the mean-field reference with the self-consistent distortion, the proper density response must be the time-dependent mean field in the space of electron states at $\vect{k}$ and $\vect{k} + \vect{Q}_{\mathrm{CDW}}$. Its kernel is the functional derivative of the same self-energy that produced $\Delta_{\mathrm{CDW}}$: Hartree--Fock exchange plus the interaction mediated by the phonon at $\vect{Q}_{\mathrm{CDW}}$. Only then does the broken translation produce the phason pole at zero frequency in the proper response, and only then does the $f$-sum rule \cref{eq:4.3-fsum} hold. The bare bubble on the reconstructed bands is the counterpart of the nonconserving bubble of \cref{sec:7.2}: it misses the phason and violates the $f$-sum. The phonon at $\vect{Q}_{\mathrm{CDW}}$ is dressed by the electronic polarization, and the interaction it mediates may enter either the electronic kernel or the dressed phonon propagator, not both. Screening with $v_{\mathrm{C}}(q)$ as in \cref{eq:15.1-Lcons} then gives the loss function, in which the longitudinal phason sits at the condensate plasma frequency. The phason and the amplitudon are composites of electrons and lattice, and assigning them to a sector of \cref{sec:8.6} is a convention, as for the pair composites of a superconductor.

A calculation for a real crystal needs that kernel on a first-principles mean field. The supercell of the distorted crystal must be relaxed with a functional that produces the distortion, and the electron--phonon vertex at $\vect{Q}_{\mathrm{CDW}}$ taken from the undistorted cell. The kernel must come from the same functional as the distortion, since a mismatch violates the Ward identity at the order of the mismatch. The checks are the Goldstone limit of the proper response, in which the phason frequency vanishes with its wave vector when $\vect{Q}_{\mathrm{CDW}}$ is made artificially incommensurate and stays gapped for a commensurate wave; the $f$-sum rule; and the measured modes, the \SI{0.951}{meV} phason of (TaSe$_4$)$_2$I among them.
